\documentclass[11pt]{article}
\usepackage[T1]{fontenc}
\usepackage[utf8]{inputenc}
\usepackage{lmodern}
\usepackage{amsmath,amssymb,amsthm}
\usepackage[margin=1in]{geometry}
\usepackage{graphicx,booktabs}
\usepackage[hidelinks]{hyperref}
\usepackage{microtype}
\hypersetup{pdftitle={Coherence time in biological oscillator assemblies bounds the rate of state registration},
pdfauthor={Ian Todd},pdfsubject={2026 author edition}}
\newtheorem{proposition}{Proposition}
\newcommand{\wrap}{\operatorname{wrap}}
\newcommand{\E}{\mathbb{E}}
\newcommand{\Prob}{\mathbb{P}}
\title{Coherence time in biological oscillator assemblies\\
bounds the rate of state registration}
\author{Ian Todd\\
Sydney Medical School, University of Sydney\\
Sydney, NSW, Australia\\
\texttt{itod2305@uni.sydney.edu.au}}
\date{2026 author edition}
\begin{document}
\maketitle

\begin{abstract}
The coherence time of a biological oscillator assembly sets the pace at
which coordinated activity becomes a registered state. This creates a
speed--flexibility trade-off: independently variable phase relations expand
the assembly's repertoire while increasing the coordination needed for a
joint readout. We derive a bound on registration rate from the probability
of satisfying the readout's requirements together, and express coherence
time as a first-passage problem in relative-phase space. The resulting
framework connects neural timing to target accessibility, opportunities
for alignment and the survival of alignment through readout. Existing
Kuramoto simulations show a strong fitted relationship in modular networks
and a weak one in sparse networks, exposing the role of architecture.
An illustrative neural calculation gives a 39 ms characteristic wait;
narrowing the phase window raises it to 163 ms. Structured phase dynamics
between registrations provide a further computational resource: they can
shape which outcome becomes available at the next readout. The framework
links the speed of biological processing to the organisation of its
internal dynamics and yields direct tests through registration survival
curves, phase perturbations and measures of temporal experience.

\end{abstract}

\noindent\textbf{Keywords:} biological oscillations; coherence time; coordination
depth; high-dimensional dynamics; state registration; temporal perception

\medskip
\noindent{\small This author edition accompanies the published article in
\textit{BioSystems} 263, 105755 (2026).
The published article remains the
\href{https://doi.org/10.1016/j.biosystems.2026.105755}
{version of record} and the article to cite.}

\section{Biological processing has a coordination timescale}

Coherence time bounds how fast a biological oscillator assembly can
register a new state. When a readout requires joint coordination, the
assembly must reach a usable configuration and sustain it through the
readout. Increasing the freedom of its phase relationships expands the
available repertoire but makes a specified joint configuration harder to
reach. This speed--flexibility trade-off connects the organisation of
biological dynamics to the pace of neural processing.

A visual scene recruits populations that respond to different features and
occupy different cortical areas. Their activity must influence downstream
circuits at useful times. Oscillatory phase affects when a population is
excitable and when an arriving signal is effective. This is the biological
motivation for communication through coherence
\cite{fries2005,fries2015,womelsdorf2007}. Working memory also involves
temporally organised bursts and interactions across frequency bands
\cite{miller2018,lundqvist2016,lisman2013}. A circuit's processing speed
therefore depends on the timing of relations among populations as well as on
the speed of its individual components.

We call a state registration a \emph{commit}: a transition that makes a
distributed configuration available as an output with sufficient persistence
to affect subsequent activity. Examples include a stable population readout,
a coordinated volley that crosses a downstream threshold, or entry into a
state maintained by recurrent activity.

The high-dimensional view of biological intelligence
\cite{todd2025intelligence} places the accessible repertoire at the centre
of computation. Here it also explains a timing cost: a readout requiring
several relationships within restricted ranges selects a small target in
that repertoire. The geometry of this target and the dynamics of reaching
it determine the waiting time.

The activity between registrations remains computationally important.
Structured phase relationships can retain recent inputs, shape
excitability and bias the next outcome. We develop the registration bound,
examine its oscillator evidence, and use it to connect continuous internal
dynamics with biological timing and temporal experience.

\section{From distributed phases to registered states}
\subsection{The configuration and the target}

Let $M$ be the number of participating modules, each an assembly with a
meaningful aggregate phase. For module $m$, containing $n_m$ oscillators
with phases $\theta_j(t)$ at time $t$, define
\begin{equation}
 z_m(t)=\frac{1}{n_m}\sum_{j\in m}e^{i\theta_j(t)}
       =r_m(t)e^{i\psi_m(t)}.
 \label{eq:module}
\end{equation}
Here $i$ is the imaginary unit, $r_m$ is within-module coherence between
zero and one, and $\psi_m$ is the aggregate phase. A phase estimate becomes
unreliable when $r_m$ is too small. Module selection and the minimum
acceptable coherence must therefore precede analysis of inter-module
alignment.

For internally coherent modules, the inter-module order parameter is
\begin{equation}
 z_{\rm inter}(t)=\frac{1}{M}\sum_{m=1}^{M}e^{i\psi_m(t)}
                =r_{\rm inter}(t)e^{i\bar\psi(t)}.
 \label{eq:inter}
\end{equation}
Its magnitude $r_{\rm inter}$ measures instantaneous phase concentration;
$\bar\psi$ is the circular mean phase. It differs from the order parameter
computed by pooling all individual oscillators.

Choose a within-module threshold $r_{\rm thr}$, a full phase tolerance
$\varepsilon$ in radians, and a dwell duration $D$ in seconds. The
alignment target $A_\varepsilon$ consists of configurations for which every
$r_m\geq r_{\rm thr}$ and
\begin{equation}
 \max_m |\wrap(\psi_m-\bar\psi)|\leq\varepsilon/2.
 \label{eq:target}
\end{equation}
The function $\wrap$ maps a phase difference into $(-\pi,\pi]$.
Configurations with undefined circular mean are outside the target.
A qualifying episode starts at a time $t$ for which the target condition
holds throughout $[t,t+D]$; registration is completed at $t+D$.
Repeated samples during the same episode count once. A new event requires
departure and rearming, or a separately specified reset.

For biological registration, the target and dwell must be linked to the
receiving circuit's response. Preferred nonzero phase lags and
cross-frequency relationships define alternative targets. The
maximum-deviation criterion checks every module; high inter-module
coherence alone can conceal an outlier.

The phase configuration lives on a product of circles. Removing the
common phase leaves $M-1$ relative-phase coordinates. This quotient
geometry supplies $M-1$ coordinates; the dynamics determine their
statistical dependence. For a smooth,
positive density near complete alignment, a sufficiently small target
has probability proportional to $\varepsilon^{M-1}$ at fixed $M$.
Confinement to a lower-dimensional support can change that exponent.
Preferred offsets, target geometry and inaccessible regions can change
the result further.

\subsection{A rate bound from necessary coordination}
\label{sec:rate}

To connect geometry to time, an \emph{attempt} must represent a distinct
opportunity for registration. Its definition includes which state starts
the opportunity and when another opportunity becomes available. It may
be a reset cycle, an independently gated sample, or a separated episode
of exploration. The attempt rate counts physical opportunities
independently of the recording frequency.

Suppose an attempt requires $k$ coordination conditions. Let $p_j$ be an
upper bound on the probability of condition $j$ being satisfied,
conditional on the earlier conditions and on the relevant attempt
history. The chain rule gives a bound on joint alignment,
$p_A\leq\prod_{j=1}^{k}p_j$, where $p_A$ denotes alignment probability
at an attempt. Let $q$ bound the conditional probability that the
alignment lasts for $D$ and meets any remaining registration requirement.
If attempts occur at asymptotic mean rate $\lambda_a$ per second, then
\begin{equation}
 \boxed{\quad
 \Gamma_{\rm reg}\leq\lambda_a q\prod_{j=1}^{k}p_j
 \leq\lambda_a q\,\bar p^{\,k},\qquad p_j\leq\bar p<1 .
 \quad}
 \label{eq:bound}
\end{equation}
Here $\Gamma_{\rm reg}$ is the long-run mean registration rate, and
$\bar p$ is a common upper bound on conditional success. The bound
assumes at most one counted registration per attempt and finite mean
attempt rate. Summing the conditional expected successes over attempts
establishes it even for dependent attempts. The probability bounds apply
uniformly over the histories being counted.

Equation~\ref{eq:bound} expresses the core coordination constraint.
Adding requirements whose conditional chances remain below one makes
successful registration exponentially scarce. Coupling can make a later
condition almost certain once earlier ones are met; directed control can
steer the state towards the target; hierarchy can replace one large
joint requirement with several smaller ones. These mechanisms reduce the
coordination cost. The number of acceptable outputs also
matters: registration into any of many targets can be faster than entry
into one specified target. Comparisons across $M$ hold the
readout task and its allowed outcomes comparable.

For an explicit waiting-time model, assume that eligible attempts have
alignment probability $p_A$ and persistence probability
$q=\Prob(\text{persistence for }D\mid\text{alignment at an attempt})$.
Eligibility here includes the within-module coherence requirement; its
availability is part of the attempt process.

\begin{proposition}[Independent attempt model]
Suppose complete attempts occupy equal intervals of duration $\Delta$,
including any required dwell and reset, and have independent success
probability $s=p_Aq>0$. Starting immediately before a new attempt, the time
$T$ to the first completed success satisfies
\begin{equation}
 \Prob(T>n\Delta)=(1-s)^n,\qquad
 \E T=\frac{\Delta}{s},\qquad n=0,1,2,\ldots.
 \label{eq:geometric}
\end{equation}
\end{proposition}
\noindent
The number of attempts is geometric: the first $n$ must all fail for the
wait to exceed $n\Delta$. Summing this survival probability proves the
mean. With $\lambda_a=1/\Delta$, the mean is
$(\lambda_a p_Aq)^{-1}$. The same mean follows for independent,
identically distributed complete renewal cycles with finite mean
duration $\Delta$, even when duration and success within a cycle are
associated; the survival law then need not be geometric.

For rare successes with weak dependence between opportunity clusters,
a constant-hazard approximation gives, with conditional success rate
$h$ per second,
\begin{equation}
 S(t)=\Prob(T>t)\simeq e^{-ht},\qquad
 \tau_{\rm coh}=h^{-1},\qquad
 t_u=-\frac{\log(1-u)}{h}.
 \label{eq:survival}
\end{equation}
Here $\tau_{\rm coh}$ is the characteristic mean waiting time and $t_u$
is the time by which a fraction $u$ of events has occurred. The median
is $\tau_{\rm coh}\log 2$ and the 95th percentile is approximately
$3\tau_{\rm coh}$. Short waits retain nonzero probability: the mean
is not a hard minimum. The positive dwell supplies a minimum time from eligible
alignment onset to completed readout. Where the model describes the
start of a successful episode, that dwell must be added to its timestamp.

\subsection{Dependence, entry and persistence}

Coherence time is a target-specific waiting time. Stationary occupancy
combines entry frequency and residence duration: rare, long visits and
frequent, short visits can produce the same occupancy. Full-process
mixing also measures something different. A slowly relaxing variable
unrelated to the target can lengthen mixing while leaving entry unchanged.
Recurrence identities for stationary discrete processes \cite{kac1947}
apply to their specified recurrence experiment; continuous-time first
entrance requires its own initial state and entry dynamics.

Measure $\lambda_a$ for the target and opportunity definition in use.
An inverse autocorrelation time of $r_{\rm inter}$ provides a proxy after
calibration against actual opportunities, including their clustering.
Its scalar summary discards phase directions. Angular frequency spread
and carrier frequency supply further timescales; neither is generally
equal to the target's attempt rate.

One empirical alternative is to count separated entries into the target
and measure the fraction that survive the dwell. If $\nu_A$ is the
long-run entry rate and $q_{\rm ent}$ the entry-conditioned survival
fraction, then $\nu_Aq_{\rm ent}$ is the rate of qualifying episodes,
subject to the rearming rule. The first-passage distribution additionally
depends on preparation: arrival from an experimental reset, recurrence
after a commit and forward waiting from a randomly chosen time have
different initial conditions and need not be exponential.

Persistence can contribute as much to scaling as entry. For example,
if each of $M$ aligned modules independently exits at rate $\beta$,
then
\begin{equation}
 q(M,D)=e^{-M\beta D}.
 \label{eq:dwell}
\end{equation}
The exit rate $\beta$ has units of inverse seconds. In this example
$-\log q=M\beta D$, so dwell adds a term growing linearly with $M$
to the logarithm of the mean wait. Use a constant persistence factor
over a range where measured survival remains approximately constant.

For a deterministic covariate schedule, or an exogenous one conditional
on its path, the survival expression is
$S(t\mid\mathcal H)=\exp[-\int_0^t h(v\mid\mathcal H)\,dv]$, where
$\mathcal H$ denotes that history and $v$ is elapsed time. Pooling
heterogeneous histories can broaden the waiting distribution. Binning
events by coherence measured at success selects unusually aligned
states; estimate pre-event hazards from the full risk set. Endogenous
phase histories require a joint model of state and event dynamics. Correlated
attempts, metastability and slow modulation should be compared using
the full survival data, including trials with no observed registration.

\section{Alignment probability and high-dimensional flexibility}
\subsection{A tractable phase-distribution model}

For a single relative-phase coordinate $\phi$, a useful closure is the
von Mises density
\begin{equation}
 f_\kappa(\phi)=\frac{e^{\kappa\cos\phi}}{2\pi I_0(\kappa)},\qquad
 p_1(\varepsilon,\kappa)=
 \int_{-\varepsilon/2}^{\varepsilon/2}f_\kappa(\phi)\,d\phi.
 \label{eq:mises}
\end{equation}
The concentration $\kappa\geq0$ controls the spread around zero;
$I_0$ and $I_1$ are modified Bessel functions of the first kind.
For this distribution, its population mean resultant length $r$
satisfies $r=I_1(\kappa)/I_0(\kappa)$ \cite{mardia2000}. The von
Mises density maximises entropy given the first circular moment.
For the closure, use the resultant length of the specified relative-phase
distribution. Instantaneous $r_{\rm inter}$ and its trajectory average are
different statistics. Substituting trajectory-average coherence for $r$
requires a distributional check using all relevant trajectories.

For $M-1$ independent anchored relative phases, each with density
$f_\kappa$, the rectangular alignment target has exact probability
$p_A=p_1^{M-1}$. The mean-centred target in
Eq.~\ref{eq:target} has a different shape and dependent deviations.
Its finite-window probability needs a geometry factor that can vary
with $M$, even where the local exponent is $M-1$.

For a declared range of tolerances and network conditions, consider
the empirical approximation
\begin{equation}
 p_A\simeq C_M p_1^{d},\qquad
 d=\alpha(M-1),\qquad
 \tau_{\rm coh}\simeq
 \frac{p_1^{-\alpha(M-1)}}{\lambda_a C_M q}.
 \label{eq:closure}
\end{equation}
The factor $C_M$ accounts for target geometry and residual distributional
effects; $d$ is a fitted probability exponent and $\alpha$ its value
per relative-phase coordinate. Require the resulting probability to lie
between zero and one. The range $0<\alpha\leq1$ describes dependencies
that facilitate joint success relative to the product reference; general
dependencies allow other exponents. Estimate $C_M$ across module counts.
The empirical closure and the history-uniform bounds in
Eq.~\ref{eq:bound} have separate evidential requirements.

At fixed $\kappa$ and small $\varepsilon$,
$p_1\simeq\varepsilon e^\kappa/(2\pi I_0(\kappa))$. Thus a constant
$d$ gives an inverse-power dependence on tolerance. Across $M$, the
logarithmic form is
\begin{equation}
 \log\tau_{\rm coh}\simeq
 -\log(\lambda_a C_Mq)
 +\alpha(M-1)[-\log p_1].
 \label{eq:log}
\end{equation}
An exponential dependence on module count is predicted when
$\alpha$ and $p_1$ stay approximately fixed and the remaining factors
do not compensate exponentially.

\subsection{Dimension depends on the question and the scale}

The directly relevant local dimension of the target probability is
the tolerance slope
\begin{equation}
 d_A(\varepsilon)=\frac{d\log p_A(\varepsilon)}
                       {d\log\varepsilon}.
 \label{eq:targetdim}
\end{equation}
This quantity measures how rapidly acceptable configurations become
scarcer as the required precision increases, holding the preparation
and attempt rule fixed. At finite tolerance it can vary with
$\varepsilon$. It agrees with $d$ in Eq.~\ref{eq:closure} only where
the closure, its constant factors and its small-window approximation
are adequate. The geometry of the support controls the asymptotic
exponent when a regular local density exists.

Covariance gives another useful view. For a specified real coordinate
representation of the relative phases, let $C$ be its covariance and
$\lambda_\ell$ its nonnegative eigenvalues. The participation ratio is
\begin{equation}
 D_{\rm PR}=\frac{(\operatorname{tr} C)^2}
                  {\operatorname{tr}(C^2)}
            =\frac{(\sum_\ell\lambda_\ell)^2}
                   {\sum_\ell\lambda_\ell^2}.
 \label{eq:pr}
\end{equation}
It measures how widely variance is distributed among modes and lies
between one and the covariance rank when variance is nonzero.
This summarises the distribution of neural activity across modes
\cite{cunningham2014,stringer2019}.

For a nonsingular Gaussian in $k$ dimensions, a sufficiently small ball
centred at a point of positive density has probability proportional to
its radius to the power $k$, however unequal the covariance
eigenvalues. Anisotropy changes its coefficient and the scales over
which directions become resolvable. Participation ratio consequently
neither supplies nor generally bounds the target exponent.
It can be tested as an empirical predictor of $d_A$ over a stated
tolerance range. Noise can increase participation ratio as well as
structured activity.

For circular data, report the representation and scale: $M-1$ anchored
phase differences, a local unwrapped chart, or a sine--cosine embedding. Wrapped deviations
from the circular mean do not generally sum to zero as real numbers.
Determine covariance rank from the chosen coordinates.

\subsection{The speed--flexibility trade-off}

The model gives a physical reason for local fast processing and slower
distributed integration to coexist. If each additional independent
phase relation can occupy several distinguishable ranges, adding
relations expands the available configuration repertoire. Requiring
one particular combination makes joint success less likely. The same
freedom that permits more possible combinations can therefore delay
a particular registration.

Within the von Mises closure, increasing concentration raises
$p_1$ and speeds alignment at fixed attempt rate, dwell and target.
Tight phase locking also suppresses variation in relative phases.
At finite noise, concentration can reduce the accessible repertoire
at the readout's resolution without changing the mathematical rank
of the support. Amplitude variation, within-module dynamics and
cross-frequency relations can remain rich. High coherence is thus
compatible with substantial dynamical complexity outside the
particular phase coordinates constrained by the readout.

Architecture controls this trade-off. Coupling changes the attempt
process, persistence and target accessibility as well as concentration.
Hierarchical organisation preserves local degrees of freedom while
reducing the number that must coordinate simultaneously. It therefore
provides a route to flexible processing with fast readout.

\section{Oscillator evidence for the coordination cost}
\subsection{Networks, detection and fitted scaling}

The existing numerical study uses noisy Kuramoto oscillators
\cite{kuramoto1984,acebron2005}. With $\omega_i$ the natural angular
frequency of oscillator $i$, $K_{ij}$ its coupling weights and
$\sigma$ the noise amplitude, the dynamics have the form
\begin{equation}
 d\theta_i=\left[\omega_i+
       \sum_j K_{ij}\sin(\theta_j-\theta_i)\right]dt
       +\sigma\,dW_i .
 \label{eq:kuramoto}
\end{equation}
The $W_i$ are independent standard Wiener processes. Coupling
normalisation is included in $K_{ij}$. The saved comparison uses
$N=100$ oscillators, module counts $M=3,\ldots,10$, 20 trials per
count, a 60 s observation horizon and a 5 ms integration step.
The natural-frequency standard deviation is 0.3 rad s$^{-1}$ and
the numerical noise amplitude is 0.1 rad s$^{-1/2}$.
All three constructions use overall coupling 0.7 s$^{-1}$ after
row normalisation. Before normalisation, modular within-module
edges have weight 1.0 and between-module edges weight 0.08.
The sparse connection probability is 0.03.

The detector requires every within-module coherence to exceed 0.6
and every module phase to lie within 0.5 radians of the circular
mean, giving full tolerance $\varepsilon=1.0$ rad. It requires six
consecutive qualifying samples, implementing a nominal 30 ms dwell
at the 5 ms step. The saved event time is the start of that run.
These are therefore waiting times to successful alignment onsets
from random initial phases. Completed readout includes the dwell;
stationary intervals require repeated events. The detector verifies
persistence at the sampled integration points.

Initial phases and natural frequencies are drawn afresh across trials;
sparse adjacency matrices are also redrawn. Event detection starts immediately;
trajectory-average coherence excludes the first quarter of each
trajectory. Non-hits are right-censored at 60 s. The fits use finite
Kaplan--Meier medians, so not every module count contributes.
Table~\ref{tab:topology} reports the saved fit summaries.

\begin{table}[tb]
\centering
\caption{Existing topology comparison. $n$ is the number of module-count
points with usable medians in the regression; each count has 20 trials.
$R^2$ describes the fitted logarithmic response, not held-out
survival prediction. The all-to-all exponent comes from a different,
linearised predictor and is not directly comparable to the other two.}
\label{tab:topology}
\begin{tabular}{llrrr}
\toprule
Topology & Predictor & $\hat\alpha$ & $R^2$ & $n$\\
\midrule
Modular & Exact $p_1$ closure & 0.989 & 0.975 & 5\\
All-to-all & Linearised surrogate & 1.137 & 0.779 & 8\\
Sparse random & Exact $p_1$ closure & 0.121 & 0.279 & 4\\
\bottomrule
\end{tabular}
\end{table}

For the exact-$p_1$ regressions, the predictor is
$(M-1)[-\log p_1(\varepsilon,\kappa(\bar r_{\rm inter}))]$,
where $\bar r_{\rm inter}$ is trajectory-average coherence pooled
over all trials. The response is the logarithm of the product of the median
waiting time and the autocorrelation-based attempt proxy.
An intercept is fitted. For all-to-all coupling, the mean window
probability is about 0.998, leaving little variation in the exact
predictor; the reported surrogate instead uses
$(1-\bar r_{\rm inter})(M-1)\log(2\pi/\varepsilon)$.

The modular network shows a strong fitted coordination cost
($R^2=0.975$). The sparse network's weak fit ($R^2=0.279$) exposes the
importance of architecture. All-to-all coupling explores a
near-synchronised regime in which alignment is readily accessible.
Together, these results place partially coordinated modular dynamics at
the centre of the present empirical case.

The saved fits use an attempt proxy averaged over hit trials and finite
medians from the less censored conditions. These selections can affect
the slope. The free intercept combines the mean--median conversion,
target geometry and persistence. Independent estimates of those factors
would turn the fitted relationship into a survival prediction.

At fixed total $N$, increasing $M$ also reduces module size, bringing
finite-size fluctuations and effective coupling into the comparison.
The accepted study reports within-module mean coherence between 0.66 and
0.74 without a declining trend, and less than 5\% change in median waits
on halving or doubling the nominal dwell. These results demonstrate
robustness over the tested range; varying module count at fixed module
size would isolate the effect of coordination depth.

The existing exponentiality diagnostic groups observed hits by
trajectory-average coherence. Its small groups describe the observed
subset. The next test uses the full censored survival curves and predicts
outcomes from a stated initialisation.

\subsection{An illustrative neural timescale}

Neural timing studies place perceptual integration on a scale of tens of
milliseconds \cite{samaha2015,cecere2015,vanrullen2016}. The phase model
reaches this scale with a compact illustrative parameter set. Behavioural
thresholds additionally include stimulus processing, decision requirements
and readout persistence.

Consider the reduced closure with $C_Mq=1$, $M=8$,
$\alpha=0.7$, $r=0.75$, full tolerance $\varepsilon=\pi/2$ and
attempt proxy $\lambda_a=126$ s$^{-1}$. These values define a scenario
rather than an independently measured human assembly. Inverting the
Bessel-function ratio and integrating Eq.~\ref{eq:mises} gives

\begin{equation}
 \kappa=2.369301,\qquad p_1=0.723078,\qquad
 \tau_{\rm coh}=\frac{p_1^{-4.9}}{126}
              =0.038871\ {\rm s}.
 \label{eq:example}
\end{equation}
The product approximation produces a characteristic wait of 39 ms.
Its alignment probability is about 0.204, outside a deep rare-event
regime, so an exponential survival law needs testing. This neural
scenario is separate from the seconds-scale numerical network sweep.

\begin{table}[tb]
\centering
\caption{Sensitivity of the illustrative closure. All rows use
$\alpha=0.7$ and $C_Mq=1$. Probabilities use numerical inversion and
quadrature, without a small-window approximation. The entries are
characteristic waiting times; any additional readout delay is separate.}
\label{tab:sensitivity}
\begin{tabular}{lrrrrr}
\toprule
Scenario & $M$ & $r$ & $\varepsilon$ &
$\lambda_a$ (s$^{-1}$) & $\tau_{\rm coh}$ (ms)\\
\midrule
Base & 8 & 0.75 & $\pi/2$ & 126 & 39\\
Fewer modules & 6 & 0.75 & $\pi/2$ & 126 & 25\\
More modules & 10 & 0.75 & $\pi/2$ & 126 & 61\\
Lower concentration & 8 & 0.65 & $\pi/2$ & 126 & 75\\
Higher concentration & 8 & 0.85 & $\pi/2$ & 126 & 19\\
Narrower window & 8 & 0.75 & $\pi/3$ & 126 & 163\\
Faster attempts & 8 & 0.75 & $\pi/2$ & 189 & 26\\
\bottomrule
\end{tabular}
\end{table}

Narrowing the window to $\pi/3$ gives $p_1=0.539752$ and
$\tau_{\rm coh}=162.883$ ms. Joint variation is broader still:
at fixed $\alpha=0.7$ and $\varepsilon=\pi/2$, the combinations
$(M,r,\lambda_a)=(6,0.85,189)$ and $(10,0.65,63)$ give 9.768
and 282.686 ms, respectively. These scenario endpoints show how
architecture and coordination shift the timescale; confidence limits
require a joint uncertainty analysis. Measuring the attempt process,
target geometry and dwell survival converts this calculation into a
prediction for a particular assembly.

\section{Predicting registration from architecture}

The central experimental prediction links architecture, precision and
timing. Estimate the coordination model on one dataset and predict the
survival curve on another. A single behavioural waiting time leaves
$\lambda_a$, $C_M$, $q$ and $\alpha$ underdetermined.

Specify the biological output and detector: module membership, frequency
bands, within-module threshold, phase tolerance, dwell and rearming.
Calibrate these against the receiving circuit's response in separate data.
A cross-correlogram half-width can inform a tolerance: temporal half-width
$w$ at angular frequency $\omega_0$ corresponds to full phase width
$2\omega_0 w$. The receiving circuit determines which width matters.
Control volume conduction and shared inputs when measuring coordination.

Estimate alignment probabilities across tolerances using eligible
opportunities from all trials. Measure target-entry rate, persistence and
temporal clustering separately, and test the attempt proxy across
conditions. Measure participation ratio and target exponent independently;
train any mapping between them before predicting the behavioural window.

Predict the held-out survival curve, including non-hits. Retain censored
trials in the likelihood and compare constant-hazard, renewal and
state-dependent-hazard models under the experimental initial state.
Propagate joint uncertainty in module selection, distribution parameters,
attempt rate and persistence, preserving their covariance.

Varying $M$ at approximately fixed module size and within-module
organisation isolates the coordination cost. Varying tolerance at fixed
opportunity process isolates the precision cost. Both manipulations predict
slower registration as the joint requirement becomes harder to satisfy.
Successful survival prediction across architectures would establish the
explanatory reach of the framework; systematic failure after independent
parameter measurement would count against it.

\section{Physical limits and the cost of a readable outcome}

Coordination places a macroscopic timing constraint on assemblies whose
individual components operate much faster. An event's latency must also
exceed every necessary physical stage: sequential stages add, while
concurrent stages overlap.

For an isolated quantum system evolving to an orthogonal state
under a time-independent Hamiltonian, the Mandelstam--Tamm and
Margolus--Levitin bounds take the form
\begin{equation}
 \tau_\perp\geq
 \max\left(\frac{\pi\hbar}{2\Delta E},
           \frac{\pi\hbar}{2(\langle E\rangle-E_0)}\right).
 \label{eq:qsl}
\end{equation}
Here $\hbar$ is the reduced Planck constant, $\Delta E$ the energy
standard deviation, $\langle E\rangle$ the mean energy and $E_0$
the ground-state energy \cite{mandelstam1945,margolus1998}.
Choosing both energy scales near $k_{\rm B}T_b$, with Boltzmann
constant $k_{\rm B}$ and temperature $T_b=310$ K, gives roughly
$4\times10^{-14}$ s. This illustrative microscopic floor lies far below
millisecond coordination. A neural assembly's open-system speed limit
instead depends on its energy distribution and interaction with its
environment.

Noise sets a further integration timescale. Detecting a weak signal takes
time even when populations already align; the signal model, noise
spectrum and decision criterion determine that delay.

If each completed operation dissipates mean energy $E_{\rm reg}$ and
available sustained power is $P$, then
$\Gamma_{\rm reg}E_{\rm reg}\leq P$. For a memory reset to a standard
state, erasing $b$ unbiased independent bits in the usual isothermal
setting dissipates at least $b k_{\rm B}T_b\ln2$
\cite{landauer1961}. The memory states, their probabilities and the reset
operation specify this cost. Participation-ratio changes and phase
threshold crossings do not identify erased bits.

A low-dimensional readout can leave the underlying high-dimensional
state active or transfer information elsewhere. Measure the physical
readout and reset costs alongside coordination and noise to identify the
assembly's rate-limiting step.

\section{Computation between registrations}
\subsection{Structured phase offsets}

The interval before registration is a computational resource. Coupled
assemblies can maintain preferred relative phases that encode recent
inputs, shape excitability and alter which target is reached next. We
propose \emph{phase-delta dynamics} as the mechanism: structured
relative-phase evolution outside the registration target that influences
future outputs.

For modules $m$ and $n$, their phase offset is
$\Delta\psi_{mn}(t)=\wrap(\psi_m(t)-\psi_n(t))$.
A computational phase pattern carries information about future outcomes
beyond shared inputs, current amplitude and other measured state
variables, and causally influences those outcomes. Persistently locked
nonzero offsets can produce low global coherence with little exploration;
the evolving offset structure supplies the relevant evidence.

Different configurations can share the same global coherence
yet favour different outcomes. This allows a concrete test:
perturb the offset structure while matching the measured
alignment opportunity rate, dwell survival and relevant
amplitude statistics as closely as possible. A change in
outcome accuracy or identity beyond that explained by altered
registration timing would support a computational role for
the pre-registration configuration. Holding mean coherence
fixed alone is insufficient, because the target probability
can still change.

This view extends communication through coherence \cite{fries2005,fries2015}.
The relation to a receiver governs when a population communicates; the
path through phase space also shapes what becomes available at readout.
Continuous computation and communication along that path prepare the
registered outcome.

\subsection{Physical substrates and spatial extent}

Phase describes physical activity sustained by synaptic interaction,
subthreshold membrane dynamics and dendritic integration. Extracellular
fields offer a further route to continuous coordination: ephaptic effects
influence cortical neurons and cerebellar spike timing
\cite{anastassiou2011,han2018}. Field-mediated organisation has also been
proposed at broader scales \cite{pinotsis2023}. Comparing field,
synaptic and common-input contributions connects the phase description
to its biological machinery.

A spatially extended narrowband signal can project many interacting modes
onto a simple waveform. A large population moving in lockstep, by
contrast, has few independent relative-phase directions. Measuring spatial
coverage, covariance structure and target probability together distinguishes
these organisations.

Under a constant-hazard model, the probability of at least one
registration within a cycle of frequency $f$ is $1-e^{-h/f}$.
To achieve reliability $u$ per cycle requires $h\geq-f\log(1-u)$.
At fixed attempt rate and with a readout required every cycle, faster
rhythms demand a more accessible target or more persistent alignment.
This timing constraint ties frequency to the coordination task; frequency
alone supplies no universal ordering of alpha and gamma dimensionality.

\subsection{Temporal experience and dynamics between registrations}

The framework separates the pace of registered outcomes from the richness
of the dynamics between them. This distinction offers a physical route
to understanding why temporal resolution, experienced duration and
retrospective duration judgements differ.

Alpha frequency has been associated with two-flash discrimination and
crossmodal temporal integration \cite{samaha2015,cecere2015}.
For a rhythm that gates opportunities with $\lambda_a$ proportional to
frequency and fixed $p_Aq$, the prediction is $\tau_{\rm coh}\propto1/f$.
A 20\% frequency increase then reduces the characteristic wait by about
16.7\%. Entrainment studies \cite{helfrich2014} provide an experimental
route: measure concentration, persistence and decision criteria alongside
frequency to distinguish their contributions.

We propose that experienced time depends partly on structured dynamics
between registrations. Richer exploration makes more relationships
available within each interval, even at an unchanged registration rate.
This predicts that structured dynamical richness contributes to prospective
duration judgements after accounting for registration timing, attention
and memory formation.

Stress exposes the distinction between duration and resolution. In a
free-fall experiment, participants retrospectively judged a frightening
event as longer without improved temporal resolution \cite{stetson2007}.
Denser memory encoding explains this dissociation. A coordination account
predicts an additional contribution only when stress changes the relevant
architecture or waiting distribution; measuring those changes separates
the mechanisms across perceptual and motor pathways.

Pharmacological perturbations provide another test. Increased signal
diversity under psychedelic drugs \cite{schartner2017,timmermann2019}
and the entropic-brain and REBUS accounts \cite{carhart2014,carhart2019rebus}
connect altered neural organisation with altered experience. The present
proposal predicts that time dilation tracks reproducible, structured
variation beyond any change in registration rate. Compare prospective
duration judgements with retrospective reports while measuring both
variables. This distinguishes computational richness from unstructured
noise, which also raises signal diversity and participation ratio.

\section{Biological implications}

Coherence time links the speed of biological processing to the freedom of
its internal dynamics. For functional events requiring simultaneous
conditions in an identifiable state space, the same coordination argument
extends beyond neural oscillators. Spatial, mechanical and biochemical
constraints define their own targets and entry processes. This provides a
common way to investigate timing in distributed biological assemblies.

Architecture offers a way to manage the cost. Local handshakes, directed
coordination and hierarchical readouts reduce the number of requirements
that must be satisfied together. These principles apply to engineered
systems as well: an asynchronous design pays for rare joint requirements,
while a global clock supplies a shared timing structure. The organisation
of the readout determines how much coordination each output demands.

\paragraph{Scope and open questions.}
The main empirical task is to predict registration survival across
architectures from independently measured entry and persistence dynamics.
The phase-delta and temporal-experience proposals add causal questions:
which offset structures change the registered outcome, and which aspects
of the intervening dynamics change experienced duration? The perturbation
and timing experiments above make these questions testable.

The speed--flexibility trade-off gives coherence time its biological
significance. Independent phase relationships expand the repertoire;
joint readout places a timing cost on accessing that repertoire.
Coupling, targeting and hierarchy control the balance. Biological
assemblies can therefore sustain rich internal dynamics while their
coordination architecture sets the pace at which those dynamics become
usable new states.

\section*{Data and code}
The accompanying accepted source and existing simulation
code are retained alongside this edition. The calculations
in Table~\ref{tab:sensitivity} use Eq.~\ref{eq:mises} and
Eq.~\ref{eq:closure}. The topology summaries use the saved
validation outputs; no new network simulations are reported.
The project repository is
\url{https://github.com/todd866/coherence-time-biosystems}.

\section*{Declarations}
No external funding supported this work. The author has
pending intellectual property related to this work.
The author thanks Abir Igamberdiev and the anonymous
\textit{BioSystems} reviewers for their feedback on the
published article.

This edition was rewritten with AI assistance (Claude, Codex)
for drafting, mathematical checking and critical review;
the author takes responsibility for the content.

\bibliographystyle{unsrt}
\bibliography{references}
\end{document}
